%% ============================================================
\chapter{Discussion}                                                                   
\label{ch:discussion}                                                                  
                 

\section{Does parallel attention improve upon recurrent fusion?}
                                                                                        
The main motivation of this work was to test whether replacing the sequential       
ConvGRU fusion in MISRGRU~\cite{resurf2024} with a Transformer's parallel
self-attention would yield a meaningful improvement in either reconstruction quality 
or inference speed:
In aggregate RSP, TRNet with a shallow transofmer (1h1d) is competitive with
but does not consistently outperform the 20-frame MISRGRU foundational model on fast 
blinking and high SNR conditions. However the attention advantage starts to show when models 
are evaluated on worse conditions where TRNet achieved consistently higher RSP. 
But on raw latency it is slower (262\,ms vs.\ 170\,ms at $N{=}20$,
Table~\ref{tab:inference_latency}).                                                    
Therefore, the attention mechanism does not give a clear answer wether to use it over convolutional
recurrence as it depends on acquisition conditions such as blinking and signal-to-noise ratio and 
compute constraints.                                                                              
                      
In TRMISR, the cross-attention softmax normalises over all input tokens, so the magnitude of the
fused representation remains stable regardless of how many frames are presented at inference      
(Table~\ref{tab:decoder_input_norm}).
Batch MISRGRU's summative GAP layer accumulates and oversaturates which degrades the output of the network 
removing the ability of feeding more input frames at inference time.
Practical advantage of TRMISR is that the     
model trained on 8 frames can be evaluated on 20, 30 or 40 frames (if the memory allows it) 
at test time with no architectural change, whereas extending MISRGRU beyond its training length 
causes progressive reconstruction degradation.                                  
This property is valuable in practice as it means that the model can be trained on shorter training sequences 
reduce compute cost and convergence time, while the researcher still has the option of feeding more frames   
for improved quality at inference. 
                                                                                        
An additional finding reinforces this picture.                                         
Counterintuitively, a TRNet trained on 8 frames outperforms one trained on 20 frames
when both are evaluated on 20- or 30-frame inputs                                      
(Figure~\ref{fig:boxplot_rsp_20fvs8f}).
We believe this happened because the training under the   
8-frame constraint forces the attention module to assign meaningful and more discriminative    
weights to every available frame.                                                         
This suggests a practical training strategy: train on the shortest frame count
that captures the relevant blinking statistics, which at the simulated acquisition     
rates corresponds to 8 frames.
                                                                                        
Ablating the Transformer fusion across four configurations (1h1d, 2h2d, 4h3d, 8h6d)    
shows that RSP does not improve meaningfully beyond 1h1d (Figure~\ref{fig:transformer_sizes}), 
while latency grows steeply: going from 1h1d to 4h3d triples the parameter count 
but increases $N{=}8$ latency by 2.4$\times$. 
This indicates that the encoder and decoder, rather than the fusion module, are the
architectural bottleneck for reconstruction quality.                          
Future improvements should focus on the encoder and decoder blocks (e.g.\ deeper residual
blocks, deformable convolutions, preprocessing such as including average widefield) 
rather than increasing the size of the transformers.                      

\section{Practicality of Streaming MISRGRU}                                
                
While the TR-MISR vs.\ MISRGRU comparison answers the thesis's primary architectural   
question, but the streaming MISRGRU variant addresses a separate and arguably more
impactful problem: real-time super-resolution during microscope acquisitions.
The batch model produces one SR output per $N$ acquired frames; the streaming model
produces one output per frame, amortising the processing cost over the entire $N$ Frame 
acquisition.
At $N{=}20$ frames on a 512$\times$512 field of view, streaming MISRGRU takes 11.29\,ms
per SR output compared to 170.29\,ms for the batch model a $15\times$ improvement in
per-output latency.

Streaming MISRGRU RSP converges to the batch-model value by the 20th accumulated frame
(Figure~\ref{fig:lln35_vs_largergrufourier}), and continues to improve beyond that     
point as the GRU accumulates more context. This capability the batch model (MISRGRU and TRNet)
lack entirely.

Having real-time feedback of reconstruction matter a lot for live-cell imaging because 
it could allow biologists to observe the live sample as it interacts with the world and interve 
if something goes wrong.
In a post-processing workflow the
low-resolution frames are collected and reconstructed only afterwards, so failures such as
focus drift, photobleaching, an incorrect exposure, or the cell moving out of the field of
view may go unnoticed until the session is over and the data is analysed. A streaming
reconstruction overcomes these problems while the experiment is still running, allowing the
operator to refocus, reposition the sample, adjust the imaging parameters before the data
quality degrades, or abort an unproductive acquisition early.

We found that the reconstruction quality saturates at around 50 accumulated frames, which sets
the practical operating point for the streaming model. At a 100\,fps acquisition rate this
corresponds to a temporal resolution of 0.5\,s. Because the model emits one super-resolved
output for every frame it receives, the reconstruction also captures motion implicitly: a
structure that moves over the accumulation window appears blurred in the output, signalling
the presence and direction of motion to the operator.  

\section{Motion reconstruction robustness}                                                         
Real-time imaging of moving structures was achieved with a              
motion-augmented training.
We compared finetuning the baseline model with adding 
Deformable Phase-space Alignment (DPA) module.                                                                                
At 0.2\,px/frame motion speed the DPA variant achieved similar RSP score to its static reference
in the fast-high SNR condition, Finetuning the model also worked, however, during visual inspection
DPA module showed less blur artefacts.

Importantly, the motion-finetuned models does not retain as good performance on static data 
as TRMISR or Streaming MISRGRU (Figure~\ref{fig:motion_snr_static}). 

Degradation concentrated in the slow conditions where temporal averaging is most valuable.
This trade-off is expected: a model regularised to handle inter-frame displacement     
cannot simultaneously exploit the long-range temporal correlation that static,    
slow-blinking acquisitions provide.                                                    
For deployment, the static streaming model is prefered when the sample is
known to be stationary or the drift is below the model's tolerance; the                
motion-finetuned or DPA variant should be selected when the sample is expected to exhibits visible  
drift relative to the field of view.

% \section{Reflection}

% Deep learning models emerge more often 

% This work sits within a broader shift in the data-science and AI field: the replacement of
% analytic scientific pipelines with learned surrogates that approximate the same output far
% faster. Classical SOFI is a closed-form statistical estimator: it computes cumulants of the
% pixel intensity traces and inverts the imaging model explicitly. 

% The deep-learning models
% studied in this thesis discard can reach better results 

% that explicit computation and instead learn the mapping from a short
% frame sequence to the SOFI image directly. The same pattern recurs across scientific
% imaging, from learned reconstruction in MRI and cryo-EM to neural surrogates for physical
% simulation. A recurring theme in this thesis reinforces the value of cross-domain transfer
% in this trend: both fusion architectures originate in satellite multi-image
% super-resolution and were carried, with relatively small adaptations, into fluorescence
% microscopy. That a method developed for fusing satellite views of the Earth transfers to
% fusing blinking fluorophores illustrates how architectural ideas in the field generalise
% across problem domains that share an underlying structure, here, the fusion of many
% low-information observations of a single static scene.

% The most important risk that follows from this shift is hallucination. A classical SOFI
% reconstruction can be noisy or low-resolution, but it cannot invent structure that the
% cumulant statistics do not support. A learned model has no such guarantee. The results of
% this thesis make the risk concrete: under slow-blinking, low-SNR conditions the MISRGRU
% baseline reconstructs filaments that are absent from the ground truth
% (Figure~\ref{fig:model_reconstruction_comparison}), and it does so confidently, with no
% indication that the output is unreliable. For a biologist, a hallucinated microtubule is
% not a cosmetic artefact; it can be mistaken for a genuine structure and drive a wrong
% conclusion. This is precisely the regime, dim signal and sparse blinking, in which
% super-resolution is most needed and least verifiable, because no independent ground truth is
% available at acquisition time. The attention mechanism mitigates but does not eliminate the
% problem: TR-MISR degrades more gracefully than the GRU, yet it too is ultimately
% interpolating from learned priors. Any deployment of these models should therefore be paired
% with uncertainty estimates or with a classical SOFI reconstruction computed in the
% background as a sanity check, so that users can distinguish recovered structure from
% fabricated detail.

% These considerations matter to several stakeholders. For the biologists and microscopists
% who would use such a tool, the benefit, super-resolved video at the camera frame rate, is
% inseparable from a new validation burden: they must learn when to trust the reconstruction
% and when to fall back on slower but trustworthy methods. For instrument vendors and software
% developers, the streaming model's $11$\,ms-per-frame inference on a single commodity GPU
% shows that real-time super-resolution is deployable without specialised hardware, lowering
% the barrier to integrating these methods into existing microscope control software. For the
% research group that produced the SOFI-MISRGRU predecessor, this thesis both validates the
% recurrent-fusion approach and maps out where it breaks down (frame-count generalisation,
% streaming throughput, motion), providing a foundation for subsequent work.

% Finally, the applicability of these findings must be qualified honestly. All quantitative
% results were obtained on synthetic data, where the ground-truth structure is known exactly
% and the imaging conditions are controlled. This is a strength for reproducibility and for
% isolating architectural effects, but it is also a limitation: real cameras, point-spread
% functions, and biological samples will differ from the simulation in ways the models have
% never seen. Until transfer to experimental data is demonstrated, the models should be
% treated as a promising methodological contribution rather than a validated instrument, and
% in particular should not be relied upon for diagnostic or clinical decisions. The
% responsible path forward is incremental, with quantitative validation against classical SOFI
% on real acquisitions before any structure recovered only by the network is treated as real.

\section{Limitations and future directions}

First, all quantitative results are obtained on a synthetic microtubule dataset
generated with the simulation pipeline from~\cite{resurf2024}. Transfer to real
experimental data is not addressed in this thesis, and the most immediate direction for
future work is to establish whether the models can reconstruct live-cell acquisitions.

A one promising direction for adapting synthetically trained models to a specific
microscope is to fine-tune the network on a sequence of blank background frames, 
that capture only detector noise and autofluorescence. This way the model could 
quickly adapt to noise levels \cite{noise2noise2018}. 

Once the models prove capable of reconstructing live-cell data, whether through
fine-tuning or out of the box, a natural next step is to pursue general-purpose
super-resolution. Most current methods~\cite{qiao2025dpatisr,resurf2024} must be fine-tuned
for the specific sample being imaged. We believe that deep-learning models could instead
operate like SOFI, generalising across acquisitions without being retrained for a particular
point-spread function or sample structure.

A natural architectural extension is a streaming TRMISR. This thesis pairs a Streaming MISRGRU
with a batch Transformer, but never a streaming Transformer, leaving one cell of the design
space unexplored. Since attention degraded more gracefully than recurrence under the harder
blinking and low-SNR conditions, a causal attention variant with a key-value cache that
emits one output per frame could combine the reconstruction robustness of TRMISR with the
throughput of the streaming GRU.

The latency results reported here are obtained with PyTorch inference. Compiling the models
with TensorRT and applying FP16 or INT8 quantisation is the obvious next step toward the
real-time target.

A further important direction is to equip the models with uncertainty estimates. The most
serious risk of replacing a closed-form estimator with a learned one is hallucination: under
slow-blinking, low-SNR conditions the models reconstruct structure that the cumulant
statistics do not support, and do so with no indication that the output is unreliable.
Per-pixel uncertainty maps, model ensembles, or a classical SOFI reconstruction computed in
the background as a sanity check would let users distinguish recovered structure from
fabricated detail, which is essential before any structure seen only in the network output is
treated as real.

Our motion augmentation also has a structural limitation: it applies a single global
displacement to the entire field of view, translating the whole structure rather than its
individual parts. Real biological motion is rarely rigid. A model trained only on globally translated
sequences may therefore generalise poorly to such local, non-rigid deformation. Although, the model 
uses the local deformable convolution to limit this dependancy on global movement,
 but it was not tested in practice. 
Extending the augmentation to per-region or non-rigid displacement fields, and pairing it with a DPA module
that estimates spatially varying rather than global alignment, is a natural next step toward
handling realistic intracellular dynamics.

Finally, training against higher-order SOFI targets could push the spatial resolution further.
This would require simulating additional data and generating the corresponding higher-order
SOFI targets, but the decoder layers can be readily adapted to estimate higher-order cumulants.
